\section{Thay đổi yêu cầu kể từ PA2}

Mục này phân biệt các thay đổi yêu cầu với nội dung làm rõ và cải tiến thiết kế. Các mục 1--3 thay đổi hoặc quy định rõ hành vi hệ thống; mục 4 chuẩn hóa quy tắc đã có trong UC-08 của PA2; mục 5 còn chưa thống nhất và cần nhóm quyết định; mục 6 ghi nhận phạm vi thiết kế giao diện, không bổ sung yêu cầu chức năng mới.

\begin{longtable}{|p{0.07\textwidth}|p{0.40\textwidth}|p{0.25\textwidth}|p{0.20\textwidth}|}
\hline

\rowcolor{TableBlue}\tableheader{\#} & \tableheader{Nội dung thay đổi} & \tableheader{Lý do} & \tableheader{Mục bị ảnh hưởng} \\ \hline
1 &
\textbf{Chuyển sang hậu kiểm (thay đổi yêu cầu):}
Bài đăng hợp lệ được công khai ngay với trạng thái
\texttt{ACTIVE}, không cần Admin duyệt trước.
Admin chỉ xem xét bài khi bài đạt ngưỡng báo cáo.
Thay đổi này xử lý mâu thuẫn giữa tiêu chí chấp nhận
của PA2 và UC-05.
&
Đơn giản hóa quy trình đăng bài, giảm khối lượng
xử lý của Admin và thống nhất luồng với UC-05, BR-03.
&
FR-03, FR-07, AC-F02-01, AC-F05-02,
UC-05, UC-09, UC-10, BR-03;
Mục 4.2, 4.4, 6.2 của PA3.
\\ \hline
2 &
\textbf{Tự động ẩn bài bị báo cáo (thay đổi yêu cầu):}
Khi có ba người dùng khác nhau báo cáo một bài,
hệ thống chuyển bài sang trạng thái \texttt{REPORTED},
ẩn bài khỏi kết quả tìm kiếm công khai và đưa bài vào
hàng đợi kiểm duyệt của Admin. PA2 đã yêu cầu đưa bài
vào hàng đợi nhưng chưa quy định nhất quán việc ẩn bài
ngay khi đạt ngưỡng.
&
Ngăn bài có nội dung gây hại hoặc gian lận tiếp tục
hiển thị công khai trong thời gian chờ kiểm duyệt.
&
FR-07, BR-03, UC-09, UC-10;
Post, Report, ModerationService, SCR-12.
\\ \hline
3 &
\textbf{Bắt buộc cung cấp thông tin liên hệ (làm rõ):}
Khi tạo bài, PA3 yêu cầu cung cấp số điện thoại hợp lệ
gồm 10 chữ số hoặc địa chỉ email. PA2 có trường liên hệ
nhưng UC-05 chưa nêu rõ trường này là bắt buộc trong
các ngoại lệ kiểm tra dữ liệu.
&
Đảm bảo người dùng có thể liên hệ sau khi Claim được
chấp nhận, đồng thời giữ kín thông tin liên hệ trước
thời điểm đó.
&
FR-03, FR-05, BR-02, NFR-02,
UC-05, UC-08;
ContactDetails, ContactAccessService,
SCR-06.
\\ \hline

4 &
\textbf{Chuẩn hóa quy tắc duyệt Claim (không thêm hành vi mới):}
PA3 quy định rõ trong BR-04 các chuyển trạng thái Claim
\texttt{PENDING} $\rightarrow$
\texttt{ACCEPTED}/\texttt{REJECTED}.
Khi Claim được chấp nhận, Post chuyển sang
\texttt{RESOLVED}; hai bên đủ quyền có thể xem
thông tin liên hệ của nhau. Nội dung này chuẩn hóa
hành vi đã được mô tả trong UC-08 của PA2.
&
Thống nhất chuyển trạng thái Claim, phân quyền
và việc cung cấp thông tin liên hệ trong toàn bộ
thiết kế.
&
FR-05, FR-06, BR-02, BR-04 mới,
UC-07, UC-08;
ClaimService, ContactAccessService,
Post, Claim.
\\ \hline

5 &
\textbf{Kiểm tra ảnh bài đăng (chưa có quyết định):}
Mục 4.1--4.2 của PA3 quy định từ 1 đến 5 ảnh JPG/PNG,
mỗi ảnh tối đa 5 MB. Tuy nhiên, UC-05 của PA2 và
SCR-06 của PA3 cho phép không tải ảnh.
Do đó, số ảnh tối thiểu vẫn chưa thống nhất.
&
Cần thống nhất quy tắc kiểm tra để giao diện,
dịch vụ ứng dụng và cơ sở dữ liệu xử lý giống nhau.
&
FR-03, UC-05, AC-F02-02, AC-F02-03;
ItemMedia, PostService, SCR-06.
\\ \hline

6 &
\textbf{Bổ sung trạng thái giao diện và bố cục di động (cải tiến thiết kế):}
PA3 mở rộng prototype PA2 với màn hình không có kết quả,
bố cục cho thiết bị rộng 360 px, hộp thoại xác nhận Claim,
thông tin liên hệ sau khi Claim được chấp nhận,
cùng các trạng thái focus và lỗi.
Đây là cải tiến thiết kế, không phải yêu cầu chức năng mới.
&
Cải thiện khả năng sử dụng, khả năng tiếp cận
và mức độ bao phủ các luồng thay thế, luồng ngoại lệ.
&
UC-03, UC-05, UC-08,
NFR-03, NFR-04;
SCR-04b, SCR-04m, SCR-05m,
SCR-06m, SCR-09b, SCR-09c;
Mục 6.3.
\\ \hline
\end{longtable}
